% Source: https://artofproblemsolving.com/community/c3h1475765_angle_bisector_and_median
% Stable identity: aops-topic-1475765-q01
Let the sides be $n-1$, $n$, $n+1$ with $n\ge2$. The largest angle is opposite the longest side $n+1$ and the smallest angle, call it $\theta$, is opposite $n-1$; the largest angle is $2\theta$.

By the law of sines,
\[
\dfrac{n+1}{n-1}=\dfrac{\sin2\theta}{\sin\theta}=2\cos\theta .
\]
By the law of cosines for the angle $\theta$,
\[
\cos\theta=\dfrac{n^2+(n+1)^2-(n-1)^2}{2n(n+1)}=\dfrac{n^2+4n}{2n(n+1)}=\dfrac{n+4}{2(n+1)} .
\]
So $\dfrac{n+1}{n-1}=\dfrac{n+4}{n+1}$, that is $(n+1)^2=(n+4)(n-1)$, or $n^2+2n+1=n^2+3n-4$, which gives $n=5$. The sides are $4,5,6$ and the longest side is $\boxed{6}$.

Check: the angle opposite $4$ has cosine $\dfrac{25+36-16}{60}=\dfrac{3}{4}$, the angle opposite $6$ has cosine $\dfrac{16+25-36}{40}=\dfrac{1}{8}$, and indeed $2\left(\dfrac{3}{4}\right)^2-1=\dfrac{1}{8}$.

\begin{asy}
usepackage("amsmath");
size(5.5cm);
defaultpen(fontsize(10pt));
pair B=(0,0),C=(6,0),A=(3.75,sqrt(25-3.75^2));
draw(A--B--C--cycle,linewidth(0.9));
draw(arc(B,0.8,0,degrees(A-B)),red);
draw(arc(A,0.6,degrees(B-A),degrees(C-A)),blue);
dot(A);dot(B);dot(C);
label("$A$",A,N); label("$B$",B,SW); label("$C$",C,SE);
label("$6$",(B+C)/2,S); label("$5$",(A+B)/2,NW); label("$4$",(A+C)/2,NE);
label("$\theta$",B+1.15*dir(degrees(A-B)/2),red);
label("$2\theta$",A+0.95*dir((degrees(B-A)+degrees(C-A))/2),blue);
label("\mbox{}",(min(currentpicture,user=true)+max(currentpicture,user=true))/2,black); // keeps the text after the figure black when LaTeX typesets the labels inline
\end{asy}
